\section{Resumen de procesos y tiempos en GPU}

Para medir el comportamiento completo del sistema se ejecutó el pipeline sobre un vídeo de \texttt{30 s} a \texttt{25 FPS}, es decir, \texttt{750} frames. Durante la ejecución se monitorizó la GPU con \texttt{nvidia-smi}. La máquina utilizada incorporaba una \texttt{NVIDIA GeForce RTX 3080} con \texttt{10 GB} de VRAM.

En esta ejecución convivieron tres procesos principales en GPU. El primero fue el servidor de lenguaje \texttt{llama-server}, que mantenía cargado Gemma en formato GGUF. El segundo fue la propia interfaz, que mantenía activo el servicio de comentarios y la síntesis XTTS. El tercero fue el proceso de tracking, donde se ejecutaban conjuntamente detección, proyección, identificación, tracking, inferencia posicional y acciones con PathCRF.

\begin{table}[ht]
\centering
\small
\begin{tabular}{lccc}
\toprule
Proceso & Función principal & VRAM media & Pico VRAM \\
\midrule
\texttt{llama-server} & Generación de texto con Gemma & 2198 MiB & 2198 MiB \\
\texttt{interfaz/app.py} & Interfaz, servidor y XTTS & 2336 MiB & 2336 MiB \\
\texttt{track.py} & Detección, tracking, acciones y PathCRF & 4807 MiB & 4942 MiB \\
\bottomrule
\end{tabular}
\caption{Procesos residentes en GPU durante la ejecución completa del pipeline.}
\end{table}

El pico total de memoria observado fue de \texttt{9515 MiB} sobre \texttt{10240 MiB}, por lo que la ejecución llegó a ocupar aproximadamente el \texttt{93\%} de la VRAM disponible. En el instante de máximo consumo convivían los tres procesos: Gemma ocupaba \texttt{2198 MiB}, la interfaz con XTTS \texttt{2336 MiB} y el proceso de tracking \texttt{4942 MiB}. Esto deja un margen libre muy pequeño, de unos \texttt{725 MiB}, lo que explica que la GPU sea un recurso crítico cuando se intenta ejecutar todo el sistema de forma simultánea.

La utilización media de GPU fue de \texttt{16.25\%}, con picos del \texttt{100\%}. Esto indica que el cuello de botella no es únicamente la capacidad de cómputo bruto de la GPU. También intervienen fases con coste en CPU, sincronizaciones entre módulos, generación secuencial de audios, escritura de vídeo y procesos de postprocesado. La VRAM, en cambio, sí queda prácticamente saturada durante la ejecución.

El tiempo total del run fue de \texttt{486.8 s} para un vídeo de \texttt{30.0 s}. Por tanto, el sistema completo se ejecutó aproximadamente \texttt{16.2} veces más lento que el tiempo real. El vídeo con tracking se guardó antes, alrededor de los \texttt{441 s}, pero el ensamblado final con comentarios, evaluación y exportación elevó el tiempo total hasta los \texttt{486.8 s}.

\begin{table}[ht]
\centering
\small
\begin{tabular}{lcc}
\toprule
Medida & Tiempo & Relación con tiempo real \\
\midrule
Duración del vídeo & 30.0 s & 1.0x \\
Vídeo con tracking generado & $\sim$441 s & $\sim$14.7x más lento \\
Pipeline completo finalizado & 486.8 s & 16.2x más lento \\
\bottomrule
\end{tabular}
\caption{Comparación entre duración del vídeo y tiempo de procesamiento.}
\end{table}

En las fases por frame, el coste medio de la parte visual fue de \texttt{492 ms} por frame. Las fases más costosas fueron la proyección al campo y la inferencia posicional, ambas con unos \texttt{155 ms} por frame de media, seguidas de la identificación de equipo/jugador con \texttt{129 ms} por frame. La detección YOLO fue mucho más rápida en comparación, con una media de \texttt{28 ms} por frame, aunque el primer frame fue más lento por inicialización.

PathCRF se ejecutó en modo \textit{rolling} mediante checkpoints asíncronos. Se completaron \texttt{18} checkpoints y se emitieron \texttt{680} edges. El primer checkpoint tardó \texttt{18.55 s}, ya que incluyó carga del modelo desde disco y preparación inicial. Una vez cargado, los checkpoints posteriores bajaron a una media de \texttt{1.15 s}. Esta diferencia muestra la importancia de mantener los modelos residentes en memoria durante una sesión.

En cuanto a comentarios, se generaron \texttt{15} comentarios de acción más la introducción. La generación textual fue rápida: la media del LLM para acciones fue de \texttt{0.18 s}. La síntesis XTTS fue el paso más costoso dentro del audio, con una media de \texttt{4.49 s} por comentario. El primer audio de acción tardó \texttt{21.55 s}, debido al calentamiento de la voz y del backend TTS; excluyendo esa primera llamada, la media de síntesis bajó a aproximadamente \texttt{3.27 s}.

Aunque se generaron \texttt{15} audios de acción, el vídeo final solo mezcló \texttt{6} eventos. Esto se debe al modo diferido: los comentarios que se solapan temporalmente con un audio anterior no se desplazan hacia delante, sino que se descartan de la pista sonora para mantener coherencia temporal con la imagen. Por tanto, el sistema generó más comentarios de los que finalmente se escuchan.

En conjunto, la prueba muestra que el pipeline completo funciona de extremo a extremo, pero todavía está lejos del tiempo real. La RTX 3080 permite ejecutar todos los módulos, pero con la VRAM prácticamente llena y con un tiempo total unas \texttt{16} veces superior a la duración del vídeo. Para acercarse al directo sería necesario reducir el coste por frame, mantener todos los modelos precalentados, evitar regeneraciones de voz, optimizar la exportación de vídeo y probablemente introducir un pequeño retardo controlado en la visualización.